\documentclass[12pt]{article}
\usepackage{amsmath, amssymb, amsthm}
\usepackage{centernot}
\usepackage{geometry}
\geometry{margin=2.5cm}

% Simple manual environments - no amsthm numbering involved
\newenvironment{satz}[1]{%
  \par\vspace{6pt}\noindent\textbf{Theorem #1.}\itshape\space
}{%
  \par\vspace{6pt}
}

\newenvironment{definition}[1]{%
  \par\vspace{6pt}\noindent\textbf{Definition #1.}\itshape\space
}{%
  \par\vspace{6pt}
}

\newenvironment{bemerkung}[1]{%
  \par\vspace{6pt}\noindent\textbf{Remark #1.}\itshape\space
}{%
  \par\vspace{6pt}
}

\renewcommand{\qedsymbol}{$\blacksquare$}
\setcounter{secnumdepth}{0}

\title{6. Function Sequences}
\date{}

\begin{document}
\maketitle

\noindent\textbf{Convention:} $(X,\|\cdot\|_X)$ is a normed vector space and $(Y,\|\cdot\|_Y)$ is a Banach space.

\section{6.1 Uniform Convergence}

\begin{definition}{6.1.1}
Let $M \subset X$, $M \neq \emptyset$, and $f, f_n$, $n \in \mathbb{N}_0$, be mappings from $M$ to $Y$.

We say $f_n$ converges pointwise to $f$, in short
\[
f_n \xrightarrow[n\to\infty]{\mathrm{pkt}} f,
\]
if for every $x \in M$:
\[
\lim_{n\to\infty} f_n(x) = f(x).
\]

We say $f_n$ converges uniformly to $f$, in short
\[
f_n \xrightarrow[n\to\infty]{\mathrm{glm}} f,
\]
if:
\[
\forall \epsilon > 0\ \exists N_\epsilon\ \forall n \ge N_\epsilon\ \forall x \in M:
\|f_n(x)-f(x)\|_Y \le \epsilon.
\]
\end{definition}

\begin{bemerkung}{6.1.2}
\hfill
\begin{enumerate}
    \item $f_n \xrightarrow[n\to\infty]{\mathrm{pkt}} f \;\Longrightarrow\; (f_n(x))_{n\in\mathbb{N}_0}$ is a Cauchy sequence for every $x$.
    \item $f_n \xrightarrow[n\to\infty]{\mathrm{glm}} f \;\Longrightarrow\; f_n \xrightarrow[n\to\infty]{\mathrm{pkt}} f$.
\end{enumerate}
\end{bemerkung}

\par\vspace{6pt}\noindent\textbf{Example 6.1.3.}\itshape\space
\begin{enumerate}
    \item $M=[0,1]$, $f_n(x)=x^n$.
    \[
    f(x):=
    \begin{cases}
    0, & x\in[0,1[,\\
    1, & x=1.
    \end{cases}
    \]
    Then:
    \[
    f_n \xrightarrow[n\to\infty]{\mathrm{pkt}} f.
    \]

    \item $M=[0,1]$ and
    \[
    f_n(x):=
    \begin{cases}
    2nx, & 0\le x\le \frac{1}{2n},\\[1mm]
    2-2nx, & \frac{1}{2n}<x<\frac{1}{n},\\[1mm]
    0, & \frac{1}{n}\le x\le 1.
    \end{cases}
    \]
    and $f(x):=0$, $\forall x\in M$. Then:
    \[
    f_n \xrightarrow[n\to\infty]{\mathrm{pkt}} f,
    \qquad\text{but}\qquad
    f_n \not\xrightarrow[n\to\infty]{\mathrm{glm}} f.
    \]

    \item $M=\mathbb{R}_{>0}$ and
    \[
    f_n(x):=
    \begin{cases}
    1, & n\le x\le n+1,\\
    0, & \text{otherwise},
    \end{cases}
    \qquad f=0.
    \]
    Then:
    \[
    f_n \not\xrightarrow[n\to\infty]{\mathrm{glm}} f.
    \]

    \item $M=\mathbb{R}_{>0}$ and $f_n(x)=\frac{1}{nx}$. Then:
    \begin{enumerate}
        \item $f_n \xrightarrow[n\to\infty]{\mathrm{pkt}} 0$.
        \item $f_n \xrightarrow[n\to\infty]{\mathrm{glm}} 0$ on $\left[c,\infty\right[$ for all $c>0$.
        \item $f_n \not\xrightarrow[n\to\infty]{\mathrm{glm}} 0$ on $\mathbb{R}_{>0}$.
    \end{enumerate}
\end{enumerate}
\upshape

\begin{proof}
Exercise.
\end{proof}

% ...existing code...

\begin{definition}{6.1.4}
Let $B(M,Y)$ be the set of all bounded functions from $M$ to $Y$.  
For $f \in B(M,Y)$ define
\[
\|f\|_\infty := \sup_{x\in M}\|f(x)\|_Y .
\]
\end{definition}

\begin{bemerkung}{6.1.5}
\[
f_n \xrightarrow[n\to\infty]{\mathrm{glm}} f
\quad\Longleftrightarrow\quad
\|f_n-f\|_\infty \xrightarrow[n\to\infty]{} 0.
\]
\end{bemerkung}

\begin{satz}{6.1.6}
$(B(M,Y),\|\cdot\|_\infty)$ is a Banach space.
\end{satz}

\begin{proof}
As in the case $C(M,Y)$, one shows that $B(M,Y)$ is a vector space and that $\|\cdot\|_\infty$ is a norm.

Let $(f_n)$ be a Cauchy sequence in $B(M,Y)$. Then for each fixed $x\in M$, $(f_n(x))$ is a Cauchy sequence in $Y$, hence convergent (since $Y$ is Banach). Define
\[
f(x):=\lim_{n\to\infty}f_n(x).
\]

It remains to show $f\in B(M,Y)$, i.e.\ $f$ is bounded.  
Using the reverse triangle inequality, $(a_n)$ with $a_n:=\|f_n\|_\infty$ is a Cauchy sequence in $\mathbb R$, hence bounded. Let
\[
C:=\sup_{n\in\mathbb N_0} a_n.
\]
For each $x\in M$ there exists $n_x\in\mathbb N_0$ such that
\[
\|f(x)-f_{n_x}(x)\|_Y\le 1.
\]
Therefore
\[
\|f(x)\|_Y
\le
\|f_{n_x}(x)\|_Y+\|f(x)-f_{n_x}(x)\|_Y
\le C+1,
\]
for all $x$, so $f$ is bounded.

Now prove convergence: let $\varepsilon>0$. Since $(f_n)$ is Cauchy in $\|\cdot\|_\infty$, there exists $n_\varepsilon$ such that
\[
\|f_n-f_m\|_\infty<\varepsilon
\qquad \forall n,m\ge n_\varepsilon.
\]
For $x\in M$ and $n\ge n_\varepsilon$:
\[
\|f(x)-f_n(x)\|_Y
=
\left\|\lim_{m\to\infty}(f_m(x)-f_n(x))\right\|_Y
\le
\limsup_{m\to\infty}\|f_m-f_n\|_\infty
\le \varepsilon.
\]
Hence $\|f-f_n\|_\infty\le\varepsilon$ for $n\ge n_\varepsilon$, i.e.
\[
\|f-f_n\|_\infty \xrightarrow[n\to\infty]{} 0.
\]
So $B(M,Y)$ is complete.
\end{proof}

\begin{satz}{6.1.7 (Cauchy criterion for uniform convergence)}
The following statements are equivalent:
\begin{enumerate}
    \item $f_n$ converges uniformly.
    \item $\forall \epsilon>0\ \exists N_\epsilon\in\mathbb N\ \forall n,m\ge N_\epsilon:\ \|f_n-f_m\|_\infty<\epsilon$.
\end{enumerate}
\end{satz}

\begin{proof}
$(1)\Longrightarrow(2)$\\
We have $f_n \xrightarrow[n\to\infty]{\mathrm{glm}} f$ and hence
$f_n-f \xrightarrow[n\to\infty]{\mathrm{glm}} 0$.
Therefore $(f_n-f)$ is a Cauchy sequence in $B(M,Y)$. Finally,
\[
\|f_n-f_m\|_\infty\le \|f_n-f\|_\infty+\|f-f_m\|_\infty.
\]

$(2)\Longrightarrow(1)$\\
There exists $N_1$ with $\|f_n-f_{N_1}\|_\infty<1$ for all $n\ge N_1$.
Set $\hat f:=f_{N_1}$. Then $(f_n-\hat f)$ is a Cauchy sequence in $B(M,Y)$ and
therefore converges with respect to $\|\cdot\|_\infty$ to some $\tilde f\in B(M,Y)$.
Hence
\[
f_n \xrightarrow[n\to\infty]{\mathrm{glm}} \hat f+\tilde f.
\]
\end{proof}

\begin{definition}{6.1.8}
The function series $\sum f_n$ is called
\begin{enumerate}
    \item pointwise convergent,
    \item absolutely convergent,
    \item uniformly convergent,
\end{enumerate}
if:
\begin{enumerate}
    \item $\sum f_n(x)$ converges for every $x\in M$,
    \item $\sum \|f_n(x)\|_Y$ converges for every $x\in M$,
    \item the sequence of partial sums $s_n(x):=\sum f_n(x)$ converges uniformly.
\end{enumerate}
\end{definition}

\begin{bemerkung}{6.1.9}
\hfill
\begin{enumerate}
    \item $\sum f_n$ absolutely convergent $\Longrightarrow$ $\sum f_n$ pointwise convergent.
    \item $\sum f_n$ uniformly convergent $\centernot\Longrightarrow$ $\sum f_n$ absolutely convergent.
\end{enumerate}
\end{bemerkung}

\begin{proof}
\begin{enumerate}
    \item Theorem 3.5.2 (absolute convergence (for numerical series) implies convergence (for numerical series)).
    \item Choose
    \[
    f_n(x)=\frac{(-1)^n}{n}
    \qquad \forall n\in\mathbb N,\ \forall x\in M.
    \]
    Then $\sum f_n$ converges uniformly, but not absolutely.

    Let $M=(-1,1)$, $Y=\mathbb R$, $f_n(x)=x^n$. Then $\sum f_n$ is absolutely convergent.
    For uniform convergence:
    \[
    \left\|\sum_{k=n+1}^{m} f_k\right\|_\infty
    =\sup_{-1<x<1}\left|\sum_{k=n+1}^{m}x^k\right|
    =\sup_{-1<x<1}\left|x^{n+1}\frac{1-x^{m-n}}{1-x}\right|
    =\lim_{x\to1}\left(x^{n+1}\frac{1-x^{m-n}}{1-x}\right)
    =m-n,
    \]
    i.e.\ $\sum f_n$ is not uniformly convergent.
\end{enumerate}
\end{proof}

\begin{satz}{6.1.10}
Let $\sum a_k(x-x_0)^k$ be a power series in $\mathbb K$ with radius of convergence $\rho>0$.
Then $\sum a_k(x-x_0)^k$ converges for every $r\in\,]0,\rho[$ uniformly and absolutely on $\overline{B}(x_0,r)$.
\end{satz}

\begin{proof}
Let $r\in\,]0,\rho[$, $M:=\overline{B}(x_0,r)$, $Y=\mathbb K$, and
\[
f_k(x)=a_k(x-x_0)^k.
\]
Then:
\[
\|f_k\|_\infty=|a_k|\,r^k
\qquad\text{and}\qquad
\sum_k \|f_k\|_\infty \text{ converges.}
\]

(Use the definition of the radius of convergence
\[
\limsup_{k\to\infty}\sqrt[k]{\|f_k\|_\infty}
= r\limsup_{k\to\infty}\sqrt[k]{|a_k|}
= \frac r\rho <1
\]
and the root test.)

From ``$\sum \|f_k\|_\infty$ converges'' it follows that
$\sum \|f_k(x)\|_Y$ converges, i.e.\ ``$\sum f_k$ converges absolutely''.
From ``$\sum \|f_k\|_\infty$ converges'' it follows that
\[
\left\|\sum_{k=n+1}^{m} f_k\right\|_\infty
\le
\sum_{k=n+1}^{m}\|f_k\|_\infty.
\]
\end{proof}

\end{document}